% Recuperacion de 09c_entropia_sectorial.tex, 11_confluencia.tex y
% 85_agujero_negro_doble_hoja.tex; originales en fuentes_conservadas/
% horizontes_informacion. Procedencia: desarrollo/HORIZONTES_INFORMACION_FUENTES.md.
% Las referencias a antecedentes se acuerdan con el editor de VII.
% No incorpora una corriente material ni evalua la amplitud torsional.

\section{Horizontes, purificación y balances de información}
\label{vii:sec:horizontes-informacion}

El estado enriquecido conserva las rutas y sus registros durante la
evolución; una lectura de subsistema determina qué parte de esa información
queda accesible. La construcción de esta sección utiliza los canales
angulares de la sección~\ref{v:ant:sec:ii-angulos}, la incidencia y el
operador de área de la sección~\ref{v:ant:sec:area}, y las unidades
de acción, área y temperatura de las
secciones~\ref{v:ant:sec:gravity}
y~\ref{v:ant:sec:ii-kb-aut-desarrollo}. Son representaciones anteriores de la
misma construcción APP--TRIT--TPK. Cada una transmite aquí su dominio,
orientación, normalización y condiciones de realización.

Se distinguen tres espacios: el registro binario de hoja, los factores
tríticos de la incidencia sectorial y la realización geométrica de los
exteriores de un horizonte. La purificación, el transporte de observables
y la conversión dimensional precisan las relaciones entre ellos. Esta
distinción conserva la estructura que acompaña a cada valor entrópico.

\subsection{Registro de hoja y purificación bipartita}
\label{vii:subsec:hojas-purificacion}

Escribamos \(a=A_{\rm rad}\) y \(y=C^*_{\rm rad}\) para las
coordenadas angulares ya construidas. Los canales y sus probabilidades son
\begin{equation}
 \mathfrak q_\pm=e^{-a\mp y},\qquad
 p_\pm=\frac{\mathfrak q_\pm}{\mathfrak q_++\mathfrak q_-}
       =\frac{e^{\mp y}}{2\cosh y}.
 \label{vii:eq:probabilidades-hoja}
\end{equation}
Los signos indexan estos dos canales. En
\(\mathcal H_{\rm h}=\operatorname{span}\{|+\rangle,|-\rangle\}\)
definimos
\[
 \rho_{\rm h}(y)=p_+|+\rangle\langle+|+p_-|-\rangle\langle-|.
\]
El espacio auxiliar
\(\mathcal H_{\rm ar}=\operatorname{span}
\{|\mathrm{adv}\rangle,|\mathrm{ret}\rangle\}\)
porta dos etiquetas ortonormales. Fijada una fase \(\chi_0\), el estado
conjunto es
\begin{equation}
 |\Psi_{\rm ar}\rangle
 =\sqrt{p_+}|+\rangle\otimes|\mathrm{adv}\rangle
  +e^{i\chi_0}\sqrt{p_-}|-\rangle\otimes|\mathrm{ret}\rangle.
 \label{vii:eq:purificacion-hoja}
\end{equation}

\begin{proposition}[Información accesible en el registro de hoja]
\label{vii:prop:purificacion-hoja}
El vector~\eqref{vii:eq:purificacion-hoja} está normalizado y su traza
parcial sobre \(\mathcal H_{\rm ar}\) es \(\rho_{\rm h}(y)\).
Su entropía de entrelazamiento, expresada en nats, es
\begin{equation}
 s_{\rm h}(y)=\log(2\cosh y)-y\tanh y.
 \label{vii:eq:entropia-hoja}
\end{equation}
El intercambio de las etiquetas \(+\) y \(-\) conjuga
\(\rho_{\rm h}(y)\) con \(\rho_{\rm h}(-y)\), y conserva
\(s_{\rm h}\).
\end{proposition}
\begin{proof}
La norma al cuadrado es \(p_++p_-=1\). La ortogonalidad de las
etiquetas auxiliares elimina los términos cruzados de la traza parcial.
Los coeficientes de Schmidt al cuadrado son precisamente \(p_\pm\).
Por~\eqref{vii:eq:probabilidades-hoja},
\(p_+-p_-=-\tanh y\); sustituir
\(\log p_\pm=\mp y-\log(2\cosh y)\) en
\(-\sum_\pm p_\pm\log p_\pm\) demuestra la fórmula.
El intercambio de signos permuta los dos autovalores y conserva la entropía.
\end{proof}

El factor común \(e^{-a}\) se cancela al normalizar el registro binario.
La fase relativa \(\chi_0\) permanece en el estado conjunto aunque esta
reducción diagonal no la publique. La reversibilidad lógica se refiere
al transporte del registro completo, con sus etiquetas y correlaciones;
el retorno de fase conserva el avance de memoria descrito en la
sección~\ref{vii:sec:retorno-memoria-gravedad}. Una implementación térmica
tiene además su propia ley de energía y disipación. En particular,
\(s_{\rm h}\) mide este subsistema binario y su bipartición concreta.

\subsection{Purificación de los sectores de incidencia}
\label{vii:subsec:purificacion-sectorial}

La incidencia de borde construida anteriormente distingue dos conjuntos
de dieciocho enlaces, con etiquetas sectoriales \(m=90,120\). Cada enlace
porta \(\mathbb C^3\), con base \(|0\rangle,|1\rangle,|2\rangle\)
y operador de ocupación \(\operatorname{diag}(0,1,2)\). Sean
\(\mathsf N_m\) sus sumas sectoriales, y
\begin{equation}
 \mathsf N=\mathsf N_{90}+\mathsf N_{120},\qquad
 \mathsf D_\partial=90\mathsf N_{90}+120\mathsf N_{120}.
 \label{vii:eq:ocupaciones-sectoriales}
\end{equation}
En la rama positiva, \(\gamma_{\rm BI}\) denota el parámetro de
Barbero--Immirzi ya obtenido, y \(a_0>0\) su factor areal angular.
Esta notación distingue el parámetro de la conexión y de sus holonomías.
El área se expresa como
\begin{equation}
 \widehat{\mathcal A}_m=\ell_P^2\gamma_{\rm BI}a_0\mathsf N_m,
 \qquad
 \widehat{\mathcal A}=\widehat{\mathcal A}_{90}
                         +\widehat{\mathcal A}_{120}.
 \label{vii:eq:area-sectores}
\end{equation}

Para el parámetro finito \(\Delta=\Delta_{\rm mod}\in\mathbb R\) del
estado reducido, ponemos
\begin{align}
 x_m&=e^{-m\Delta},& z_m&=1+x_m+x_m^2,&
 Z(\Delta)&=z_{90}^{18}z_{120}^{18},\\
 \rho_\partial(\Delta)&=Z(\Delta)^{-1}e^{-\Delta\mathsf D_\partial}.
 \label{vii:eq:estado-sectorial}
\end{align}
La traza en esta sección es la traza matricial ordinaria y los estados
tienen traza uno. El parámetro \(\Delta\) conserva la selección y la
normalización del estado antecedente.

\begin{proposition}[Purificación tensorial de la frontera]
\label{vii:prop:purificacion-sectorial}
Para un peso real finito \(w\), sea
\begin{equation}
 |\operatorname{TFD}(w)\rangle
 =\frac{|0\rangle|0\rangle+e^{-w/2}|1\rangle|1\rangle
                   +e^{-w}|2\rangle|2\rangle}
        {\sqrt{1+e^{-w}+e^{-2w}}}.
 \label{vii:eq:tfd-local}
\end{equation}
El producto
\begin{equation}
 |\Psi_\Delta\rangle
 =|\operatorname{TFD}(90\Delta)\rangle^{\otimes18}
  \otimes|\operatorname{TFD}(120\Delta)\rangle^{\otimes18}
 \label{vii:eq:tfd-sectorial}
\end{equation}
es una purificación normalizada de \(\rho_\partial(\Delta)\).
\end{proposition}
\begin{proof}
En cada enlace, la suma de los cuadrados de los tres coeficientes es uno.
La traza parcial elimina los términos con índices distintos y devuelve
\(\operatorname{diag}(1,e^{-w},e^{-2w})/(1+e^{-w}+e^{-2w})\).
La traza parcial conmuta con el producto tensorial. Multiplicar las
treinta y seis reducciones da el exponencial y la función de partición
de~\eqref{vii:eq:estado-sectorial}.
\end{proof}

\subsection{Ecuación informacional y fluctuaciones areales}
\label{vii:subsec:ecuacion-informacional}

Las medias y varianzas de ocupación por enlace son
\begin{equation}
 f_m=\frac{x_m+2x_m^2}{z_m},\qquad
 v_m=\frac{x_m+4x_m^2}{z_m}-f_m^2.
 \label{vii:eq:momentos-sectoriales}
\end{equation}
Escribimos \(s(\rho)=-\operatorname{Tr}(\rho\log\rho)\), con
\(0\log0=0\). En la purificación anterior, \(s(\rho_\partial)\)
es su entropía de entrelazamiento.

\begin{proposition}[Área media, entropía y respuesta sectorial]
\label{vii:prop:respuesta-sectorial}
Para \(\Delta\) real y finito se cumple
\begin{align}
 \langle\mathsf N\rangle_\Delta&=18(f_{90}+f_{120}),\\
 \langle\mathsf D_\partial\rangle_\Delta&=18(90f_{90}+120f_{120}),\\
 \langle\widehat{\mathcal A}\rangle_\Delta
 &=18\ell_P^2\gamma_{\rm BI}a_0(f_{90}+f_{120}),\\
 s(\Delta)&=\log Z(\Delta)+\Delta\langle\mathsf D_\partial\rangle_\Delta,
 \label{vii:eq:entropia-sectorial}\\
 \frac{\mathrm d\langle\mathsf D_\partial\rangle_\Delta}{\mathrm d\Delta}
 &=-\operatorname{Var}_\Delta(\mathsf D_\partial)
  =-18(90^2v_{90}+120^2v_{120}),\\
 \frac{\mathrm ds}{\mathrm d\Delta}
 &=-\Delta\operatorname{Var}_\Delta(\mathsf D_\partial),\\
 \frac{\mathrm d\langle\widehat{\mathcal A}\rangle_\Delta}{\mathrm d\Delta}
 &=-18\ell_P^2\gamma_{\rm BI}a_0(90v_{90}+120v_{120}).
 \label{vii:eq:respuestas-area-entropia}
\end{align}
\end{proposition}
\begin{proof}
Los pesos locales son \((1,x_m,x_m^2)/z_m\), cuyos dos primeros
momentos dan \(f_m\) y \(v_m\). La aditividad de los operadores
ocupación y la factorización del estado dan las tres primeras fórmulas.
Tomar la traza de
\(-\rho_\partial\log\rho_\partial
=\rho_\partial(\log Z+\Delta\mathsf D_\partial)\)
da la entropía. La diferenciación de las sumas finitas produce
\(\partial_\Delta\log Z=-\langle\mathsf D_\partial\rangle_\Delta\)
y \(\partial_\Delta^2\log Z
=\operatorname{Var}_\Delta(\mathsf D_\partial)\).
Los factores independientes suman sus varianzas. En la derivada de
\eqref{vii:eq:entropia-sectorial} se cancelan los dos términos de
primer momento. Finalmente, \(\partial_\Delta f_m=-mv_m\)
da la respuesta del área.
\end{proof}

La pérdida de uniformidad se cuantifica en dimensión \(d_\partial=3^{36}\):
\begin{equation}
 I_\partial(\Delta)
 =36\log3-s(\Delta)
 =\mathcal D\bigl(\rho_\partial(\Delta)\Vert I/d_\partial\bigr)\ge0.
 \label{vii:eq:deficit-informacional}
\end{equation}
Aquí \(\mathcal D\) denota la entropía relativa, definida y justificada
abajo. Sustituir \(\log(I/d_\partial)=-36\log3\,I\) prueba la
identidad. La convexidad de \(t\log t\) da la positividad; la igualdad
corresponde a la distribución uniforme, que en esta colección ocurre
exactamente en \(\Delta=0\). El déficit mide información adimensional;
la escala \(\gamma_{\rm BI}a_0\ell_P^2\) mide área por ocupación.

\begin{theorem}[Inversión de ocupaciones y área complementaria]
\label{vii:thm:inversion-sectorial}
Sea \(\mathcal R_\partial\) la involución que intercambia
\(|0\rangle\) y \(|2\rangle\) y fija \(|1\rangle\) en cada
factor. Entonces
\begin{align}
 \mathcal R_\partial\mathsf N\mathcal R_\partial^*&=72I-\mathsf N,\\
 \mathcal R_\partial\mathsf D_\partial\mathcal R_\partial^*
 &=7560I-\mathsf D_\partial,\\
 \rho_\partial(-\Delta)&=\mathcal R_\partial\rho_\partial(\Delta)
                          \mathcal R_\partial^*,\\
 s(-\Delta)&=s(\Delta),\\
 \langle\widehat{\mathcal A}\rangle_{-\Delta}
 +\langle\widehat{\mathcal A}\rangle_\Delta
 &=72\gamma_{\rm BI}a_0\ell_P^2,\\
 \mathcal D\bigl(\rho_\partial(\Delta)\Vert\rho_\partial(-\Delta)\bigr)
 &=\Delta\bigl(7560-2\langle\mathsf D_\partial\rangle_\Delta\bigr).
 \label{vii:eq:inversion-sectorial}
\end{align}
\end{theorem}
\begin{proof}
Cada ocupación se transforma por \(n\mapsto2-n\). En cada sector,
\(\mathsf N_m\mapsto36I-\mathsf N_m\), por lo que el término
constante del lector ponderado es \(36(90+120)=7560\).
Además, \(z_m(-\Delta)=e^{2m\Delta}z_m(\Delta)\) y
\(Z(-\Delta)=e^{7560\Delta}Z(\Delta)\). La sustitución en el
estado normalizado prueba la conjugación. Su espectro conserva la entropía;
la transformación de \(\mathsf N\) da el área complementaria. Restar
los logaritmos de los dos estados y tomar la esperanza en
\(\rho_\partial(\Delta)\) da la última igualdad.
\end{proof}

En \(\Delta=0\), \(s=36\log3\). Cuando \(\Delta\to+\infty\),
la ocupación nula concentra el estado y la entropía tiende a cero;
la inversión lleva ese límite a la ocupación máxima. Para \(\Delta>0\),
la entropía decrece estrictamente porque el lector ponderado tiene más
de un autovalor y el estado tiene rango completo. La configuración
completa, el área y la entropía conservan así funciones diferenciadas.

\subsection{Ley modular y variaciones finitas de área}
\label{vii:subsec:ley-modular-finita}

Fijemos \(\rho_0=\rho_\partial(\Delta_0)\), con
\(\Delta_0\) real finito, y mantengamos
\(\Delta_0,\gamma_{\rm BI},a_0,\ell_P\) constantes. Las áreas
normalizadas \(\mathcal A_m^{\rm n}=\gamma_{\rm BI}a_0\mathsf N_m\)
expresan el operador modular como
\begin{equation}
 \mathsf K_0=-\log\rho_0
 =c_0I+\sum_{m=90,120}\beta_m\mathcal A_m^{\rm n},
 \quad \beta_m=\frac{m\Delta_0}{\gamma_{\rm BI}a_0},
 \quad c_0=\log Z(\Delta_0).
 \label{vii:eq:operador-modular-areal}
\end{equation}
Se tiene \(\beta_{120}=\tfrac43\beta_{90}\), también cuando ambos
son cero. El cociente es \(4/3\) cuando \(\Delta_0\ne0\).
\(\mathsf K_0\) es el logaritmo de este estado de referencia; su tipo
lo distingue del transporte temporal y de la conexión gravitatoria.

\begin{theorem}[Primera ley modular y corrección finita]
\label{vii:thm:ley-modular-finita}
Para cualquier estado normalizado \(\sigma\) del mismo espacio, sea
\(\Delta_\sigma\langle B\rangle
=\operatorname{Tr}[(\sigma-\rho_0)B]\). Entonces
\begin{equation}
 s(\sigma)-s(\rho_0)
 =\sum_{m=90,120}\beta_m\Delta_\sigma\langle\mathcal A_m^{\rm n}\rangle
  -\mathcal D(\sigma\Vert\rho_0),
 \label{vii:eq:ley-modular-finita}
\end{equation}
donde
\(\mathcal D(\sigma\Vert\rho_0)
=\operatorname{Tr}[\sigma(\log\sigma-\log\rho_0)]\)
es finita y no negativa. Si \(\sigma(\varepsilon)\) es diferenciable,
normalizada y \(\sigma(0)=\rho_0\), se obtiene
\begin{equation}
 \left.\frac{\mathrm ds(\sigma(\varepsilon))}{\mathrm d\varepsilon}\right|_0
 =\sum_{m=90,120}\beta_m
   \left.\frac{\mathrm d\langle\mathcal A_m^{\rm n}\rangle_{
                 \sigma(\varepsilon)}}{\mathrm d\varepsilon}\right|_0.
 \label{vii:eq:primera-ley-modular}
\end{equation}
\end{theorem}
\begin{proof}
Por definición,
\(\mathcal D(\sigma\Vert\rho_0)=-s(\sigma)
+\operatorname{Tr}(\sigma\mathsf K_0)\), mientras
\(s(\rho_0)=\operatorname{Tr}(\rho_0\mathsf K_0)\).
La resta y~\eqref{vii:eq:operador-modular-areal} prueban la identidad
finita, pues la normalización cancela el término \(c_0I\).

Para la positividad, diagonalicemos \(\sigma\) con autovalores
\(p_i\) y vectores \(u_i\), y \(\rho_0\) con autovalores
\(q_j>0\) y vectores \(v_j\). Si
\(r_i=\langle u_i,\rho_0u_i\rangle>0\), la concavidad del
logaritmo da
\[
 \sum_j|\langle u_i,v_j\rangle|^2\log q_j\le\log r_i.
\]
En consecuencia,
\(\mathcal D(\sigma\Vert\rho_0)
\ge\sum_i p_i\log(p_i/r_i)\ge0\).
La última desigualdad es la convexidad de \(t\log t\) con pesos
\(r_i\), cuya suma es uno. La fidelidad de \(\rho_0\) mantiene
finitos sus logaritmos; los autovalores nulos de \(\sigma\) se
interpretan mediante \(0\log0=0\).

Para la fórmula infinitesimal, en el entorno del estado fiel \(\rho_0\),
la derivada de la traza de \(-\sigma\log\sigma\) es
\(-\operatorname{Tr}[\dot\sigma(\log\rho_0+I)]\).
Esta identidad se obtiene de la derivada espectral de la traza.
La condición \(\operatorname{Tr}\dot\sigma=0\) deja
\(\operatorname{Tr}(\dot\sigma\mathsf K_0)\), y sustituir
\eqref{vii:eq:operador-modular-areal} concluye la prueba.
\end{proof}

Para las áreas físicas, los coeficientes son \(\beta_m/\ell_P^2\).
La aplicación del transductor común \(k_B\) da \(S=k_Bs\) y
conserva íntegro el término \(k_B\mathcal D\). La derivada en
\eqref{vii:eq:respuestas-area-entropia} varía \(\Delta\); la ley
\eqref{vii:eq:primera-ley-modular} mantiene fijo \(\Delta_0\) y
varía el estado alrededor de la referencia. Son operaciones distintas.

\subsection{Acción, energía de decisión y unidad areal}
\label{vii:subsec:conversion-informacion-area}

Fijamos una misma sección orientada de acción y sus secciones gravitatoria
y térmica, en la rama positiva. La composición
de~\ref{g26:sec:rigidez-radial-es} identifica
\(\ell_P=L_\alpha=L\) en esa realización. En la sección de retorno,
la energía característica es \(E_L=\hbar_{\rm ret}c/L\). Se tiene
\begin{equation}
 \begin{aligned}
 \ell_P^2&=\frac{\hbar G}{c^3},& t_P&=\frac{\ell_P}{c},\\
 T_{108}&=2\pi t_P,& h&=2\pi\hbar,\\
 E_P&=\frac{\hbar}{t_P}=k_B\Theta_{\rm clk}\log3.&&
 \end{aligned}
 \label{vii:eq:unidad-accion-informacion}
\end{equation}
La igualdad térmica conserva la pareja y la normalización previamente
construidas; al cambiar su base dimensional se transportan conjuntamente
\(k_B\) y \(\Theta_{\rm clk}\). El incremento elemental de área es
\(\delta\mathcal A=\gamma_{\rm BI}a_0\ell_P^2\).

\begin{proposition}[Coeficiente de conversión entre decisión y área]
\label{vii:prop:tension-informacion-area}
El cociente de la energía de decisión y del incremento areal satisface
\begin{equation}
 \begin{aligned}
 \mathcal T_{I/A}:=\frac{E_P}{\delta\mathcal A}
 &=\frac{k_B\Theta_{\rm clk}\log3}{\gamma_{\rm BI}a_0\ell_P^2}\\
 &=\frac{c^3}{\gamma_{\rm BI}a_0G t_P}
  =\frac{c^4}{\gamma_{\rm BI}a_0G\ell_P}.
 \end{aligned}
 \label{vii:eq:tension-informacion-area}
\end{equation}
Su dimensión es energía por área.
\end{proposition}
\begin{proof}
La sustitución de \(E_P=\hbar/t_P\) y
\(\ell_P^2=\hbar G/c^3\) cancela la sección de acción común.
La identidad \(\ell_P=ct_P\) da la última expresión. La positividad
de los factores asegura que todos los cocientes están definidos.
\end{proof}

La capacidad \(81\) del corte cúbico y la bipartición de los treinta y
seis enlaces sectoriales corresponden a dominios diferentes. En el
grafo cúbico \(\{0,\ldots,N-1\}^3\), las \(N^2\) rectas
paralelas entre dos caras opuestas son disjuntas por enlaces: todo corte
intercepta las \(N^2\), y una capa transversal alcanza esa cota.
Para \(N=9\), el estado producto uniforme de los trits de esa capa
tiene \(s_{\rm corte}=81\log3\). La asignación de una decisión del
mismo ciclo a cada enlace expresa
\begin{equation}
 S_{\rm corte}^{\rm fis}=81k_B\log3,\qquad
 E_{\rm corte}^{(1)}=81E_P
                  =\Theta_{\rm clk}S_{\rm corte}^{\rm fis}.
 \label{vii:eq:corte-energia-informacion}
\end{equation}
La prueba combina el corte mínimo, la aditividad de la entropía del
estado producto y~\eqref{vii:eq:unidad-accion-informacion}. En el estado
sectorial, la ecuación finita pertinente sigue siendo
\eqref{vii:eq:ley-modular-finita}, con su entropía relativa.

\subsection{Normalización cosmológica del horizonte}
\label{vii:subsec:balance-horizonte}

Para un radio areal \(R>0\), la realización cosmológica considerada
utiliza la tríada
\begin{equation}
 E_\Lambda(R)=\frac{c^4R}{2G},\qquad
 S_{\rm H}(R)=\frac{k_B\pi R^2}{\ell_P^2},\qquad
 T_{\rm cos}(R)=\frac{\hbar c}{2\pi k_BR}.
 \label{vii:eq:horizonte-cosmologico}
\end{equation}
El factor \(2\pi\) pertenece a esta normalización de temperatura.
El área \(\mathcal A_{\rm H}=4\pi R^2\) expresa la misma entropía
como \(S_{\rm H}/k_B=\mathcal A_{\rm H}/(4\ell_P^2)\).

Para conservar la diferencia entre escala construida y radio variable,
escribimos aquí \(R=R_h=\zeta_hL_\alpha\), con \(\zeta_h>0\).
En la sección de retorno, sustituir el mismo \(G\) y la misma área
de referencia da
\[
 E_\Lambda(R_h)=\frac{\zeta_h}{2}E_L,\qquad
 \frac{S_{\rm H}(R_h)}{k_B}=\pi\zeta_h^2,\qquad
 k_BT_{\rm cos}(R_h)=\frac{E_L}{2\pi\zeta_h}.
\]
Cada igualdad procede directamente de
\eqref{vii:eq:horizonte-cosmologico}. El factor \(1/2\) pertenece
a esta energía de horizonte. La lectura reducida de una masa
\(m_{P,\rm ret}y\) tiene radio \(r_g=L_\alpha y\); su radio de
Schwarzschild es \(2L_\alpha y\). La condición geométrica de un
horizonte cosmológico mantiene su dominio propio.

\begin{theorem}[Celda gravitatoria de media acción]
\label{vii:thm:celda-media-accion}
En~\eqref{vii:eq:horizonte-cosmologico},
\begin{align}
 T_{\rm cos}(R)S_{\rm H}(R)&=E_\Lambda(R),\\
 T_{\rm cos}(\ell_P)S_{\rm H}(\ell_P)
 &=\frac{E_P}{2}=\frac{k_B\Theta_{\rm clk}\log3}{2},\\
 E_\Lambda(\ell_P)t_P&=\frac{\hbar}{2},&
 E_\Lambda(\ell_P)T_{108}&=\frac h2.
 \label{vii:eq:media-accion-horizonte}
\end{align}
Manteniendo fijas las constantes de la coordenatización y variando \(R\),
\begin{equation}
 T_{\rm cos}\,\mathrm dS_{\rm H}=2\,\mathrm dE_\Lambda,
 \qquad S_{\rm H}\,\mathrm dT_{\rm cos}=-\mathrm dE_\Lambda.
 \label{vii:eq:diferencial-horizonte}
\end{equation}
\end{theorem}
\begin{proof}
Multiplicar temperatura y entropía da
\(\hbar cR/(2\ell_P^2)=c^4R/(2G)\).
En \(R=\ell_P\), el resultado es
\(\hbar/(2t_P)=E_P/2\). Las dos duraciones de
\eqref{vii:eq:unidad-accion-informacion} dan las relaciones de acción.
Para los diferenciales,
\(\mathrm dE_\Lambda=c^4\mathrm dR/(2G)\),
\(\mathrm dS_{\rm H}=2\pi k_BR\mathrm dR/\ell_P^2\) y
\(\mathrm dT_{\rm cos}=-T_{\rm cos}\mathrm dR/R\).
La sustitución prueba ambas identidades.
\end{proof}

El producto \(E_\Lambda=T_{\rm cos}S_{\rm H}\) expresa una
identidad de estado. Sus variaciones radiales conservan los dos términos
de~\eqref{vii:eq:diferencial-horizonte}; una interpretación termodinámica
completa especifica además el generador temporal, el trabajo y la
realización material correspondientes. En esta coordenatización, acción y
gravitación fijan la unidad \(\ell_P^2\), mientras Boltzmann convierte
información en entropía dimensional. El transporte de secciones con
\(G\hbar\) fijo conserva esa unidad cuando la sección \(c\) se
mantiene fija, y conserva \(S_{\rm H}/k_B\) a radio fijo.

En la celda \(R=\ell_P\), \(S_{\rm H}/k_B=\pi\), de modo que
la multiplicidad efectiva es \(\Omega_{\rm eff}=e^\pi\).
Para la involución hiperbólica trítica \(H_{\rm trit}\), con espectro
\(\{1,-1\}\), la diagonalización da
\(\operatorname{spec}(e^{\pi H_{\rm trit}})=\{e^\pi,e^{-\pi}\}\).
Esta relación espectral especifica el valor expansivo; una identificación
microscópica de ambos espacios conserva adicionalmente su mapa de
realización y el estado.

% Ampliación aditiva: acción, área, memoria sectorial y termometría.
% Propietarios: 70_horizontes_e_informacion.tex; sección térmica antecedente;
% 71_energia_libre_holografica_cerrada.tex del corpus integral.
% Se mantienen las condiciones de la realización cosmológica declarada.

\subsection{Conversión térmica entre el calendario y el horizonte}
\label{viii:subsec:conversion-termica-calendario-horizonte}

La acción, el incremento de área y la memoria angular admiten una
composición energética explícita. Se conserva una misma orientación de
acción, la rama positiva de Barbero--Immirzi y la realización de
\eqref{vii:eq:horizonte-cosmologico}. Abreviamos
\[
 L=\ell_P=L_\alpha,\qquad
 E_P=\frac{h}{108t_0}=\frac{\hbar c}{L},\qquad
 q_{\mathcal A}=\delta\mathcal A=\gamma_{\rm BI}a_0L^2.
\]
El símbolo \(q_{\mathcal A}\) denota aquí un incremento de área.
El factor angular \(a_0\), el funcional \(\gamma_{\rm BI}\) y
la acción son las magnitudes obtenidas en las secciones precedentes.
Su dominio conserva la incidencia hexádico--octádica, la orientación y
el calendario del mismo estado enriquecido APP--TRIT--TPK.

\begin{proposition}[Composición de las escalas areal y térmica]
\label{viii:prop:composicion-area-termica}
Con la tensión de~\eqref{vii:eq:tension-informacion-area}, se cumple
\begin{equation}
 \frac{\mathcal T_{I/A}q_{\mathcal A}}{\log3}
 =\frac{E_P}{\log3}
 =\frac{h}{108t_0\log3}
 =k_B\Theta_{\rm clk}.
 \label{viii:eq:composicion-area-termica}
\end{equation}
Para \(R=\zeta L\), con \(\zeta>0\), sea
\(\tau_R=k_BT_{\rm cos}(R)\). Entonces
\begin{equation}
 \tau_R=\frac{E_P}{2\pi\zeta},\qquad
 \frac{T_{\rm cos}(R)}{\Theta_{\rm clk}}
 =\frac{\log3}{2\pi\zeta}.
 \label{viii:eq:conversion-dos-temperaturas}
\end{equation}
En particular, el cociente en \(R=L\) es \(\log3/(2\pi)\).
\end{proposition}
\begin{proof}
La definición \(\mathcal T_{I/A}=E_P/q_{\mathcal A}\) y la
normalización trítica de~\eqref{vii:eq:unidad-accion-informacion}
dan~\eqref{viii:eq:composicion-area-termica}. La expresión
gravitatoria de la misma tensión conduce al mismo resultado:
\[
 \frac{c^4}{\gamma_{\rm BI}a_0GL}
        \gamma_{\rm BI}a_0L^2
 =\frac{c^4L}{G}=\frac{\hbar c}{L}=E_P,
\]
donde se ha utilizado \(G\hbar=c^3L^2\).
Por~\eqref{vii:eq:horizonte-cosmologico},
\(\tau_R=\hbar c/(2\pi R)=E_P/(2\pi\zeta)\).
Dividir esta igualdad por \(k_B\Theta_{\rm clk}=E_P/\log3\)
prueba el cociente de temperaturas.
\end{proof}

En la misma sección de acción, la composición gravitatoria queda
expresada por
\begin{equation}
 Gk_B\Theta_{\rm clk}\log3=GE_P=c^4L,
 \qquad G\tau_R=\frac{c^4L}{2\pi\zeta}.
 \label{viii:eq:composicion-termica-gravitatoria}
\end{equation}
Ambas identidades siguen de \(G\hbar=c^3L^2\) y de las escalas
anteriores.

La tensión convierte el incremento areal en la energía de decisión.
La normalización trítica y la periodicidad cosmológica determinan dos
secciones térmicas de esa misma energía, relacionadas por
\eqref{viii:eq:conversion-dos-temperaturas}. La conversión conserva
ambas normalizaciones y la función de Barbero--Immirzi en la resolución
geométrica.

\subsection{Respuesta energética de la memoria sectorial}
\label{viii:subsec:respuesta-energetica-sectorial}

Fijemos \(R>0\) y el estado fiel
\(\rho_0=\rho_\partial(\Delta_0)\), con \(\Delta_0>0\),
de~\eqref{vii:eq:estado-sectorial}. Su operador modular es
\(\mathsf K_0=(\log Z_0)I+\Delta_0\mathsf D_\partial\).
La realización de horizonte asigna a cada estado de densidad
\(\sigma\) las magnitudes de respuesta
\begin{equation}
\begin{aligned}
 \mathcal E_R(\sigma)
 &=E_\Lambda(R)+\tau_R
   \operatorname{Tr}[(\sigma-\rho_0)\mathsf K_0],\\
 \mathcal S_R(\sigma)
 &=S_{\rm H}(R)+k_B[s(\sigma)-s(\rho_0)].
\end{aligned}
 \label{viii:eq:respuesta-horizonte-definida}
\end{equation}
Se trata de la respuesta del estado modular en la realización geométrica
ya especificada. El radio y la acción fijan su escala energética
\(\tau_R\); las ocupaciones conservan la procedencia de cada canal.

\Needspace{28\baselineskip}
\begin{theorem}[Hamiltoniano de respuesta y balance informacional]
\label{viii:thm:hamiltoniano-respuesta-termica}
El operador de diferencias energéticas correspondiente a
\eqref{viii:eq:respuesta-horizonte-definida} puede elegirse como
\begin{equation}
 \mathsf H_R=\tau_R\Delta_0\mathsf D_\partial.
 \label{viii:eq:hamiltoniano-respuesta}
\end{equation}
Su estado de Gibbs a escala energética \(\tau_R\) es \(\rho_0\).
Para todo estado de densidad \(\sigma\) en el espacio sectorial,
\begin{equation}
 \mathcal E_R(\sigma)-T_{\rm cos}(R)\mathcal S_R(\sigma)
 =\tau_R\mathcal D(\sigma\Vert\rho_0)\geq0.
 \label{viii:eq:balance-libre-termico}
\end{equation}
En una variación diferenciable de estados alrededor de \(\rho_0\),
con \(R\) y \(\mathsf H_R\) fijos, se obtiene
\begin{equation}
 \delta\mathcal E_R
 =\tau_R\,\delta s
 =T_{\rm cos}(R)\,\delta\mathcal S_R.
 \label{viii:eq:primera-ley-energetica}
\end{equation}
\end{theorem}
\begin{proof}
Las trazas de \(\sigma\) y \(\rho_0\) son uno, por lo que la
diferencia elimina el término escalar \(\log Z_0\). De aquí resulta
\(\mathcal E_R(\sigma)-E_\Lambda(R)
=\operatorname{Tr}[(\sigma-\rho_0)\mathsf H_R]\).
Además,
\[
 \frac{e^{-\mathsf H_R/\tau_R}}
 {\operatorname{Tr}e^{-\mathsf H_R/\tau_R}}
 =\frac{e^{-\Delta_0\mathsf D_\partial}}{Z_0}=\rho_0.
\]
La igualdad \(E_\Lambda=T_{\rm cos}S_{\rm H}\) cancela la
referencia geométrica en el balance. La identidad finita
\eqref{vii:eq:ley-modular-finita}, escrita como
\(s(\sigma)-s(\rho_0)
=\operatorname{Tr}[(\sigma-\rho_0)\mathsf K_0]
-\mathcal D(\sigma\Vert\rho_0)\), demuestra
\eqref{viii:eq:balance-libre-termico}. Su término tangente es la
primera ley modular ya probada y produce
\eqref{viii:eq:primera-ley-energetica}.
\end{proof}

Para un incremento unitario de ocupación, los saltos energéticos son
\(\epsilon_{90}=90\tau_R\Delta_0\) y
\(\epsilon_{120}=120\tau_R\Delta_0\). En ambos sectores,
\begin{equation}
 \frac{\epsilon_m}{m\Delta_0}=\tau_R,
 \qquad m\in\{90,120\}.
 \label{viii:eq:pendiente-termica-comun}
\end{equation}
Así, el cociente \(4/3\) entre pesos areales registra los diferentes
saltos del mismo Hamiltoniano. El cociente entre un salto energético y
la variación correspondiente de \(-\log p\) conserva una escala común.
La normalización del estado elimina sólo una constante escalar.

\begin{corollary}[Susceptibilidad energética a Hamiltoniano fijo]
\label{viii:cor:respuesta-escala-termica}
Para \(\rho_\tau=e^{-\mathsf H_R/\tau}/
\operatorname{Tr}e^{-\mathsf H_R/\tau}\), con \(\tau>0\),
\begin{equation}
 \frac{\mathrm d\langle\mathsf H_R\rangle_\tau}{\mathrm d\tau}
 =\frac{\operatorname{Var}_{\rho_\tau}(\mathsf H_R)}{\tau^2}>0.
 \label{viii:eq:susceptibilidad-termica}
\end{equation}
En \(\tau=\tau_R\), el resultado es
\(\Delta_0^2\operatorname{Var}_{\rho_0}(\mathsf D_\partial)\).
\end{corollary}
\begin{proof}
La derivación de la suma finita de exponenciales da la varianza dividida
por \(\tau^2\). El estado es fiel y \(\mathsf H_R\) tiene más de
un autovalor, de modo que la varianza es positiva. La sustitución de
\eqref{viii:eq:hamiltoniano-respuesta} da la evaluación indicada.
\end{proof}

Esta derivada varía la escala de Gibbs y mantiene fijo el Hamiltoniano.
La variación radial de~\eqref{vii:eq:diferencial-horizonte} modifica
también \(\mathsf H_R\) y conserva sus términos geométricos propios.

\subsection{Sección termométrica y conservación del registro}
\label{viii:subsec:seccion-termometrica-registro}

El estado \(\rho_0\) y el Hamiltoniano de respuesta permiten fijar una
coordenada termométrica anterior a cualquier lector numérico candidato.
Para el radio fijado, sea \(\mathcal L_\Theta=\mathbb R u_R\) una
recta orientada cuya sección positiva \(u_R\) corresponde al estado
de referencia. En los estados fieles de Gibbs del mismo Hamiltoniano,
\begin{equation}
 \rho_\beta=\frac{e^{-\beta\mathsf H_R}}
                  {\operatorname{Tr}e^{-\beta\mathsf H_R}},\qquad
 \Theta_R(\rho_\beta)=\frac{1}{\beta\tau_R}u_R,
 \quad \beta>0.
 \label{viii:eq:lector-termometrico}
\end{equation}
La identidad operatoria
\(-\log\rho_\beta=\beta\mathsf H_R+(\log Z_\beta)I\)
determina \(\beta\): restar dos autovalores con distinta energía
da su cociente único. En particular,
\(\Theta_R(\rho_0)=u_R\).
La recta \(\mathcal L_E\) es el espacio unidimensional de las
magnitudes energéticas.

\Needspace{23\baselineskip}
\begin{theorem}[Transductor positivo y criterio de identificación]
\label{viii:thm:identificacion-transductor}
Existe un único mapa lineal positivo
\(B_R:\mathcal L_\Theta\to\mathcal L_E\) con
\(B_R(u_R)=\tau_R\). Satisface
\begin{equation}
 B_R(\Theta_R(\rho_\beta))=\beta^{-1}.
 \label{viii:eq:transductor-gibbs}
\end{equation}
La sección cronológica compatible es
\begin{equation}
 \Theta_{{\rm clk},R}
 =\frac{E_P}{\tau_R\log3}u_R
 =\frac{2\pi R}{L\log3}u_R,
 \qquad
 B_R(\Theta_{{\rm clk},R})=
 \frac{\mathcal T_{I/A}q_{\mathcal A}}{\log3}.
 \label{viii:eq:seccion-cronologica-compatible}
\end{equation}
Para cualquier otro mapa lineal \(B_*\) entre las mismas rectas,
\begin{equation}
 B_*=B_R
 \quad\Longleftrightarrow\quad
 B_*(\Theta_{{\rm clk},R})=\frac{h}{108t_0\log3}.
 \label{viii:eq:criterio-una-seccion}
\end{equation}
\end{theorem}
\begin{proof}
La sección \(u_R\) es una base de una recta, de modo que prescribir
su imagen define un único mapa lineal; la positividad de \(\tau_R\)
asegura la positividad del mapa. Sustituir
\eqref{viii:eq:lector-termometrico} prueba
\eqref{viii:eq:transductor-gibbs}. La expresión de \(\tau_R\) y
\eqref{viii:eq:composicion-area-termica} dan
\eqref{viii:eq:seccion-cronologica-compatible}. Finalmente,
\(\Theta_{{\rm clk},R}\) es no nula. Dos mapas lineales de una
recta coinciden si y sólo si coinciden sobre esa sección.
\end{proof}

El criterio compara lectores en una sección obtenida previamente del
estado, la respuesta energética y la normalización trítica. La aplicación
a una coordenada numérica concreta requiere evaluar ese lector en la
sección especificada. Las bases dimensionales se transportan conjuntamente
con los mapas y con las coordenadas térmicas que expresan.

\begin{proposition}[Transporte de la respuesta con memoria]
\label{viii:prop:transporte-termico-memoria}
Sea \(J\) una isometría que conserva las etiquetas completas del
registro, y considérese su imagen como espacio receptor. Para
\(\mathsf H'_R=J\mathsf H_RJ^*\), \(\rho'_0=J\rho_0J^*\)
y \(\sigma'=J\sigma J^*\), se conservan las expectativas energéticas,
las entropías y el balance~\eqref{viii:eq:balance-libre-termico}.
También se cumple \(J\mathsf H_R=\mathsf H'_RJ\).
\end{proposition}
\begin{proof}
Sobre la imagen, \(J\) es unitario y transporta el cálculo funcional.
La ciclicidad de la traza conserva expectativas; los mismos autovalores
no nulos conservan entropía y entropía relativa. La identidad de
entrelazamiento usa \(J^*J=I\). Sustituir estas igualdades en el balance
prueba su conservación.
\end{proof}

La misma equivalencia admite otra escala \(\tau'=a\tau_R\),
con \(a>0\), y otro origen energético \(e'\). Si
\(\mathsf H'=e'I+aJ\mathsf H_RJ^*\), entonces
\[
 \frac{\mathsf H'-e'I}{\tau'}
 =J\frac{\mathsf H_R}{\tau_R}J^*,\qquad
 \frac{e^{-\mathsf H'/\tau'}}
      {\operatorname{Tr}_{\operatorname{im}J}e^{-\mathsf H'/\tau'}}
 =J\rho_0J^*.
\]
El factor escalar \(e^{-e'/\tau'}\) se cancela en el estado;
las diferencias energéticas se multiplican por \(a\), y la entropía
relativa permanece invariante. El balance correspondiente tiene, por
tanto, coeficiente \(\tau'\). La expectativa energética completa
conserva el origen: \(\langle\mathsf H'\rangle
=e'+a\langle\mathsf H_R\rangle\). La comparación de realizaciones
se expresa así mediante sus exponentes centrados y sus escalas propias.
Si la expectativa alimenta las lecturas \(m=E/c^2\) y
\(r_g=GE/c^4\), mantener \(c,G\) da
\[
 m'=\frac{e'}{c^2}+am,\qquad
 r'_g=\frac{Ge'}{c^4}+ar_g.
\]
La normalización estadística y el transporte de magnitudes mantienen,
por tanto, funciones distintas para el mismo origen registrado.

La composición enlaza el incremento areal de Barbero--Immirzi, la energía
del calendario y la distribución angular del borde. Su transporte
retiene la procedencia sectorial y el archivo: el retorno de fase
continúa acompañado por el avance de memoria del estado enriquecido.


\subsection{Coordenatización de Kruskal e intercambio de los exteriores}
\label{vii:subsec:kruskal-dos-hojas}

Fijemos un estado expresado \(x\), su radio \(R=R(x)>0\) y las
secciones positivas \(c,G\). En la coordenatización que sigue, \(x,R,c,G\)
son parámetros fijos; \(t,r\) son las coordenadas espacio-temporales.
La realización exterior de las dos hojas se especifica por
\begin{equation}
 f(r)=1-\frac Rr,\qquad
 \mathrm ds^2=-f(r)c^2\mathrm dt^2+f(r)^{-1}\mathrm dr^2
                 +r^2\mathrm d\Omega_2^2,\qquad r>R.
 \label{vii:eq:metrica-exterior}
\end{equation}
Para el lector de hoja \(\sigma\in\{+1,-1\}\), definimos
\begin{align}
 r_*&=r+R\log\left(\frac rR-1\right),\\
 \mathcal U&=-\sigma\exp\left(-\frac{ct-r_*}{2R}\right),&
 \mathcal V&=\sigma\exp\left(\frac{ct+r_*}{2R}\right).
 \label{vii:eq:carta-kruskal}
\end{align}
Los símbolos \(\mathcal U,\mathcal V\) son coordenadas y se
distinguen de los operadores de transporte de rutas.

\begin{theorem}[Realización geométrica e involución de hoja]
\label{vii:thm:kruskal-hojas}
La coordenatización~\eqref{vii:eq:carta-kruskal} realiza las dos hojas como los
dos exteriores \(\mathcal U\mathcal V<0\) de la extensión de
Kruskal. Se cumple
\begin{equation}
 -\mathcal U\mathcal V=\left(\frac rR-1\right)e^{r/R},\qquad
 \mathrm ds^2=-\frac{4R^3}{r}e^{-r/R}\,
                    \mathrm d\mathcal U\,\mathrm d\mathcal V
                    +r^2\mathrm d\Omega_2^2.
 \label{vii:eq:metrica-kruskal}
\end{equation}
La involución
\begin{equation}
 \mathcal C_{\rm H}:(\sigma,\mathcal U,\mathcal V)
 \longmapsto(-\sigma,-\mathcal U,-\mathcal V)
 \label{vii:eq:involucion-horizonte}
\end{equation}
es isométrica, fija el radio areal, intercambia los exteriores y preserva
el conjunto de horizontes \(\mathcal U\mathcal V=0\).
La métrica es regular en \(r=R\) y satisface las ecuaciones de vacío
en el dominio regular \(r>0\).
\end{theorem}
\begin{proof}
Multiplicar las coordenadas da la primera identidad. Como
\(\mathrm dr_*/\mathrm dr=f^{-1}\),
\[
 \mathrm d\mathcal U=\frac{\mathcal U}{2R}
            (-c\,\mathrm dt+f^{-1}\mathrm dr),\qquad
 \mathrm d\mathcal V=\frac{\mathcal V}{2R}
            ( c\,\mathrm dt+f^{-1}\mathrm dr).
\]
Su producto reproduce la parte radial y temporal de la métrica.
La función \(z\mapsto(z-1)e^z\) tiene derivada \(ze^z>0\)
para \(z>0\), lo que determina \(r\) a partir del producto.
En el exterior se recuperan
\(\sigma=\operatorname{sign}\mathcal V\) y
\(t=(R/c)\log(-\mathcal V/\mathcal U)\).
En \(r=R\), el coeficiente radial de la métrica extendida es
\(-4R^2/e\), finito y distinto de cero. La inversión conserva el
producto, el radio y \(\mathrm d\mathcal U\mathrm d\mathcal V\),
e intercambia las dos elecciones de signo.

Para comprobar el vacío en la coordenatización exterior, las componentes mixtas
independientes del tensor de Einstein se reducen a
\[
 \mathsf G^t{}_t=\mathsf G^r{}_r
 =\frac{rf'+f-1}{r^2},\qquad
 \mathsf G^\theta{}_\theta=\mathsf G^\phi{}_\phi
 =\frac{f'}r+\frac{f''}2.
\]
Con \(f'=R/r^2\) y \(f''=-2R/r^3\), todas se anulan.
La expresión regular~\eqref{vii:eq:metrica-kruskal} prolonga el
resultado a través de los horizontes en el dominio \(r>0\).
\end{proof}

Esta realización conserva
\(E_{\rm H}=c^4R/(2G)\) y
\(S_{\rm H}=k_B\pi R^2/\ell_P^2\), que dependen del radio areal.
La temperatura \(T_{\rm cos}\) de
\eqref{vii:eq:horizonte-cosmologico} pertenece a la realización
cosmológica allí declarada. La coordenatización exterior
\eqref{vii:eq:metrica-exterior} conserva su propia normalización temporal.
El mapa geométrico establece el intercambio de exteriores; la comparación
con formación por colapso, espectros de radiación o modos cuasinormales
utiliza además sus condiciones físicas respectivas.

\subsection{Transporte de estados, observables y refinamientos}
\label{vii:subsec:transporte-informacion-horizonte}

Denotemos por \(F_x\) la coordenatización de los dos exteriores recién construida.
En cada coordenatización angular regular, su jacobiano temporal y radial es
\begin{equation}
 \left|\frac{\partial(\mathcal U,\mathcal V)}{\partial(t,r)}\right|
 =\frac{cr}{2R^3}e^{r/R}>0.
 \label{vii:eq:jacobiano-kruskal}
\end{equation}
La identidad resulta de los dos diferenciales de la prueba anterior.
Los restantes registros del estado se transportan como coordenadas de
fibra; la coordenatización conserva su identidad y su memoria.

\begin{proposition}[Equivalencia unitaria de las lecturas de hoja y exterior]
\label{vii:prop:transporte-unitario-horizonte}
Sea \(\mu\) una medida de hoja y \(\nu=(F_x)_*\mu\) su medida
transportada. El cambio de variables define una unitaria
\begin{equation}
 W_{\rm H}:L^2(\mu)\longrightarrow L^2(\nu),\qquad
 (W_{\rm H}\psi)(z)=\psi(F_x^{-1}z).
 \label{vii:eq:unitaria-horizonte}
\end{equation}
Con densidades respecto de medidas coordenadas fijadas, la misma
aplicación incorpora el factor de jacobiano correspondiente.
La conjugación
\begin{equation}
 \mathfrak I_{\rm H}(B)=W_{\rm H}BW_{\rm H}^*
 \label{vii:eq:mapa-observables-horizonte}
\end{equation}
es un \(*\)-isomorfismo unital completamente positivo. Transporta el
estado, sus esperanzas y su cálculo funcional modular. Para la medida
simétrica de las dos hojas, entrelaza también sus involuciones.
\end{proposition}
\begin{proof}
La definición de medida imagen da
\(\int|\psi\circ F_x^{-1}|^2\,\mathrm d\nu
=\int|\psi|^2\,\mathrm d\mu\).
La coordenatización es invertible en los exteriores, de modo que la isometría es
sobreyectiva. La conjugación unitaria preserva producto, adjunto e
identidad; en cada ampliación matricial conjuga por
\(I_n\otimes W_{\rm H}\), lo que prueba la positividad completa.
Para un estado \(\rho\), la ciclicidad de la traza da
\(\operatorname{Tr}(W_{\rm H}\rho W_{\rm H}^*
\,\mathfrak I_{\rm H}(B))=\operatorname{Tr}(\rho B)\).
El cálculo funcional implica
\(g(W_{\rm H}\rho W_{\rm H}^*)=W_{\rm H}g(\rho)W_{\rm H}^*\)
en su dominio, en particular para el logaritmo y el flujo modular de
un estado fiel en su soporte. La igualdad
\(F_x\circ\mathcal C_{\rm hoja}=\mathcal C_{\rm H}\circ F_x\)
prueba el entrelazamiento de involuciones; la medida simétrica las realiza
unitariamente.
\end{proof}

Si los refinamientos conservan \((\sigma,t,r,R)\) y utilizan la misma
isometría sobre los registros de fibra, las inclusiones
\(j_m^{\rm hoja}\) y \(j_m^{\rm ext}\) satisfacen
\begin{equation}
 W_{{\rm H},m+1}j_m^{\rm hoja}
 =j_m^{\rm ext}W_{{\rm H},m}.
 \label{vii:eq:naturalidad-horizonte}
\end{equation}
En efecto, el cambio de variables actúa sobre las coordenadas geométricas
conservadas y la inclusión sobre las mismas coordenadas de fibra; evaluar
ambas composiciones da la misma función y la misma densidad transportada.
La igualdad extiende el mapa al sistema refinado sin reiniciar los registros.

Para una bipartición \(\mathcal H_A\otimes\mathcal H_B\), un
transporte factorizado \(U_A\otimes U_B\) cumple
\[
 \rho'_A=U_A\rho_AU_A^*,\qquad s(\rho'_A)=s(\rho_A).
\]
La primera identidad sigue de la traza parcial y la segunda del espectro.
Para una unitaria general, la formulación equivalente transporta también
el álgebra del subsistema:
\(\mathcal B(\mathcal H_A)\otimes I\mapsto
U(\mathcal B(\mathcal H_A)\otimes I)U^*\).
Así se conserva la misma pregunta operacional acerca de la información
accesible. Las igualdades de entropía se aplican a estados con las entropías
finitas pertinentes.

La relación entre información y gravitación queda expresada mediante
operaciones distintas y composables: el estado y la bipartición fijan la
información accesible; la incidencia y Barbero--Immirzi fijan la resolución
areal; acción y gravitación construyen la unidad \(\ell_P^2\);
Boltzmann realiza dimensionalmente la entropía; la coordenatización de horizonte
transporta la geometría y sus observables. Las leyes finitas conservan la
entropía relativa, y las realizaciones geométricas conservan sus medidas,
su normalización temporal y sus dominios.
